\chapter{Cochlear Implant}

\section*{What Is This Surgery?}
A cochlear implant is a sophisticated, surgically implanted electronic device designed to restore a sensation of hearing to individuals with severe-to-profound sensorineural hearing loss. This intricate neuroprosthesis bypasses the damaged hair cells of the inner ear (cochlea) and directly stimulates the auditory nerve, converting sound into electrical signals that the brain interprets as sound. The procedure is typically performed under general anesthesia, involving a mastoidectomy and meticulous insertion of an electrode array into the scala tympani of the cochlea. It represents a significant advancement in otologic surgery, offering a pathway to auditory perception and speech understanding for those who derive little or no benefit from conventional amplification. The surgery aims to provide access to sound for communication, education, and improved quality of life.

\section*{Alternative Names}
\begin{itemize}
  \item Cochlear implantation
  \item Bionic ear surgery
  \item Implantable hearing device procedure
  \item Auditory prosthesis implantation
  \item Electronic ear implant surgery
  \item CI surgery
\end{itemize}

\section*{Relevant Surgical Anatomy}
Cochlear implant surgery primarily involves the temporal bone, a complex structure housing the organs of hearing and balance. Key anatomical landmarks and structures include the mastoid bone, which is drilled during mastoidectomy to create a cavity for access. Within the mastoid, the facial nerve (cranial nerve VII) courses through the fallopian canal, representing a critical structure to identify and protect due to its proximity to the surgical field and its control over facial muscle movement. The chorda tympani nerve, a branch of the facial nerve, also traverses the middle ear cavity and is often encountered during the facial recess approach.

The middle ear cavity contains the ossicles (malleus, incus, stapes) and provides access to the inner ear via the oval window (covered by the stapes footplate) and the round window. The round window membrane is a primary entry point for the electrode array into the cochlea. The cochlea itself is a snail-shaped, fluid-filled structure within the inner ear, divided into three scalae: scala vestibuli, scala media, and scala tympani. The electrode array is meticulously inserted into the scala tympani, which spirals around the modiolus, the central bony core containing the spiral ganglion neurons that form the auditory nerve. Preservation of the integrity of the facial nerve and careful manipulation around the delicate inner ear structures are paramount for a successful outcome.

\begin{itemize}
  \item \textbf{Temporal Bone:} Houses the mastoid, middle ear, and inner ear structures.
  \item \textbf{Mastoid Bone:} Air-filled cavity posterior to the ear, drilled during surgery.
  \item \textbf{Facial Nerve (CN VII):} Courses through the mastoid, critical to protect.
  \item \textbf{Chorda Tympani Nerve:} Branch of facial nerve, crosses middle ear, affects taste.
  \item \textbf{Middle Ear:} Contains ossicles, provides access to inner ear.
  \item \textbf{Round Window:} Membrane-covered opening to the scala tympani, common entry for electrodes.
  \item \textbf{Cochlea:} Snail-shaped inner ear structure, target for electrode array insertion into the scala tympani.
  \item \textbf{Auditory Nerve (CN VIII):} Receives electrical signals from the implant and transmits them to the brain.
\item \textbf{Modiolus:} Central bony core of the cochlea, containing spiral ganglion neurons.
\end{itemize}

\section*{Surgical Principle \& Rationale}
The fundamental surgical principle of a cochlear implant is to bypass the non-functional or severely damaged hair cells within the cochlea and directly stimulate the remaining viable spiral ganglion neurons of the auditory nerve. In individuals with severe-to-profound sensorineural hearing loss, these hair cells, which normally convert mechanical sound vibrations into electrical impulses, are either absent or severely impaired. Traditional hearing aids, which merely amplify sound, are ineffective in such cases because the damaged hair cells cannot transmit the amplified signals to the brain.

The rationale for cochlear implantation is to restore a functional sense of hearing by providing electrical stimulation that mimics the natural tonotopic organization of the cochlea. The electrode array, inserted into the scala tympani, delivers electrical impulses to different regions of the auditory nerve, corresponding to various sound frequencies. This direct stimulation sends signals to the brain, which learns to interpret these electrical patterns as sound. The technical goals include precise and atraumatic insertion of the electrode array, secure placement of the internal receiver-stimulator, and preservation of the facial nerve and other vital structures, all aimed at maximizing speech perception and environmental sound awareness for the recipient.

\section*{History \& Surgical Evolution}
The journey towards the modern cochlear implant began in the 1950s with pioneering experiments demonstrating that electrical stimulation of the auditory nerve could produce sound sensations. In 1957, Andr\'{e} Djourno and Charles Eyri\`{e}s in France performed the first human implant, stimulating the auditory nerve directly with a single electrode, providing rudimentary sound perception. Building on this, William House in the U.S. further advanced single-channel devices in the 1960s, proving the concept's viability and leading to the first commercially available single-channel implant in 1972.

The pivotal breakthrough came with the development of multi-channel implants, which allowed for the stimulation of different frequency regions of the cochlea, significantly improving speech understanding. Professor Graeme Clark in Australia pioneered the first successful multi-channel cochlear implant in 1978, with his patient, Rod Saunders, being the first to achieve open-set speech understanding. Independently, Ingeborg and Erwin Hochmair in Austria developed their own multi-channel device, first implanted in 1977. The U.S. Food and Drug Administration (FDA) approved multi-channel implants for adults in 1985 and for children in 1990, marking a pivotal moment in the widespread adoption and refinement of the technology. Subsequent decades have seen continuous improvements in electrode design, speech processing strategies, miniaturization, and MRI compatibility, expanding candidacy criteria and enhancing performance outcomes for a broader range of patients.

\section*{Clinical Indications}
Cochlear implantation is indicated for individuals with severe-to-profound sensorineural hearing loss who derive limited or no benefit from conventional hearing aids. Specific clinical indications include:

\begin{itemize}
  \item \textbf{Bilateral Severe-to-Profound Sensorineural Hearing Loss:} When audiometric thresholds are typically 70 dB HL or greater at 500, 1000, and 2000 Hz, and speech perception scores (e.g., sentence recognition in quiet) are below established candidacy criteria (e.g., <50\% in the best-aided condition for adults, or lack of auditory skill development in children).
  \item \textbf{Congenital or Pre-lingual Deafness in Children:} Implantation in infants and young children (as young as 9-12 months) with profound hearing loss to facilitate critical speech and language development.
  \item \textbf{Post-lingual Deafness in Adults:} For adults who acquired hearing loss after developing speech and language, and have lost the ability to understand speech despite optimal hearing aid use.
  \item \textbf{Single-Sided Deafness (SSD):} In some cases, to improve sound localization, reduce tinnitus, and enhance speech understanding in noise for individuals with profound hearing loss in one ear and normal hearing in the other.
  \item \textbf{Auditory Neuropathy Spectrum Disorder (ANSD):} For individuals with ANSD who demonstrate poor speech perception and lack of auditory development despite hearing aid use, and who meet specific audiologic and imaging criteria.
  \item \textbf{Progressive Hearing Loss:} For individuals experiencing rapidly progressive sensorineural hearing loss where hearing aids are no longer sufficient.
  \item \textbf{Meningitis-Induced Hearing Loss:} Early implantation is often recommended for hearing loss secondary to meningitis due to the risk of cochlear ossification, which can complicate later electrode insertion.
\end{itemize}

\section*{Clinical Positioning}
Cochlear implantation is positioned as a primary rehabilitative treatment for individuals with severe-to-profound sensorineural hearing loss, particularly when conventional amplification devices, such as hearing aids, have been thoroughly trialed and demonstrated to provide insufficient auditory benefit. It is considered a definitive surgical solution for restoring functional hearing in these specific patient populations.

For children with congenital or early-onset profound hearing loss, cochlear implantation is often the primary intervention to enable speech and language acquisition during critical developmental periods. In adults with post-lingual hearing loss, it serves as the primary surgical option when their ability to understand speech is severely compromised despite optimal hearing aid use. While it is a primary surgical treatment, it is always secondary to a comprehensive diagnostic evaluation of hearing loss and a trial period with appropriately fitted hearing aids, which are the initial, less invasive steps in the management pathway for most forms of hearing impairment. Thus, it is a primary surgical intervention within a broader, multi-stage rehabilitative process.

\section*{Surgical Technique \& Procedure}
Cochlear implant surgery is typically performed under general anesthesia and usually takes approximately 2-4 hours. The patient is positioned supine with the head turned to the side opposite the operative ear, and the hair around the ear is shaved or clipped. Prophylactic antibiotics are administered intravenously.

\begin{itemize}
  \item \textbf{Incision and Flap Creation:} A post-auricular incision, typically curvilinear or C-shaped, is made behind the ear. A musculoperiosteal flap is then carefully elevated to expose the mastoid bone and the squamous portion of the temporal bone.
  \item \textbf{Mastoidectomy:} A cortical mastoidectomy is performed using a high-speed surgical drill with continuous irrigation. This creates a cavity within the mastoid bone, providing access to the middle ear and the facial recess. Bone dust is meticulously cleared.
  \item \textbf{Facial Recess Approach (Posterior Tympanotomy):} A critical and delicate step involves drilling a posterior tympanotomy, or facial recess approach. This corridor is created between the facial nerve laterally and the chorda tympani nerve medially, providing a direct view into the middle ear cavity and access to the round window niche. Facial nerve monitoring is often used throughout this stage.
  \item \textbf{Cochleostomy or Round Window Insertion:} The round window membrane is identified. The surgeon either inserts the electrode array directly through the round window (round window insertion) or creates a small, precise opening (cochleostomy) anterior and inferior to the round window. The choice depends on surgical preference and anatomical considerations.
  \item \textbf{Electrode Array Insertion:} The delicate electrode array, a flexible silicone carrier with multiple platinum-iridium contacts, is carefully and slowly advanced into the scala tympani of the cochlea. The depth of insertion is crucial for optimal performance and varies by device and patient anatomy. Care is taken to minimize trauma to cochlear structures.
  \item \textbf{Receiver-Stimulator Placement:} A bone well is drilled in the temporal bone, typically superior and posterior to the mastoidectomy cavity, to securely house the internal receiver-stimulator package. The device is then fixed in place with non-absorbable sutures or small titanium screws to prevent migration.
  \item \textbf{Device Testing:} Intraoperative neural response telemetry (NRT), electrically evoked auditory brainstem response (EABR), or electrically evoked stapedial reflex threshold (ESRT) testing is performed to confirm proper device function, electrode placement, and neural response.
  \item \textbf{Closure:} The musculoperiosteal flap is repositioned, and the incision is closed in layers using absorbable sutures for the deeper tissues and non-absorbable sutures or staples for the skin. A pressure dressing may be applied to minimize swelling and hematoma formation.
\end{itemize}
No drains are typically left in place.

\section*{Preoperative Workup \& Preparation}
Thorough preoperative workup and preparation are essential to ensure patient safety, confirm candidacy, and optimize outcomes for cochlear implant surgery. This process involves a multidisciplinary team including otolaryngologists, audiologists, speech-language pathologists, and radiologists.

\begin{itemize}
  \item \textbf{Audiologic Evaluation:} Comprehensive audiometry, speech perception testing in quiet and noise, and a trial period with appropriately fitted hearing aids are performed to confirm severe-to-profound sensorineural hearing loss and insufficient benefit from conventional amplification.
  \item \textbf{Medical Clearance:} A complete medical history and physical examination are conducted. Patients undergo routine blood tests (e.g., complete blood count, coagulation profile, basic metabolic panel) and often an electrocardiogram (ECG) and chest X-ray, especially for older adults or those with significant comorbidities. Anesthesia clearance is mandatory.
  \item \textbf{Imaging Studies:} High-resolution computed tomography (CT) of the temporal bones is crucial to assess cochlear patency, identify any anatomical abnormalities (e.g., cochlear malformations, otosclerosis, ossification), and evaluate the mastoid and middle ear. Magnetic resonance imaging (MRI) of the inner ear and brain is performed to rule out retrocochlear pathology (e.g., acoustic neuroma, auditory neuropathy) and assess the integrity of the auditory nerve.
  \item \textbf{Vaccinations:} Patients, particularly children, are typically vaccinated against encapsulated bacteria, especially \textit{Streptococcus pneumoniae} and \textit{Haemophilus influenzae} type b, several weeks prior to surgery to reduce the rare but serious risk of meningitis.
  \item \textbf{Counseling and Consent:} Extensive counseling is provided to the patient and family regarding the surgical procedure, potential risks, expected outcomes, the need for intensive post-operative auditory rehabilitation, and realistic expectations for hearing improvement. Informed consent is obtained, emphasizing the commitment required for long-term success.
  \item \textbf{Medication Review:} All current medications are reviewed. Anticoagulants or antiplatelet agents (e.g., aspirin, warfarin, novel oral anticoagulants) may need to be temporarily discontinued or adjusted prior to surgery, following consultation with the prescribing physician and hematologist.
  \item \textbf{NPO Status:} Patients are instructed to be NPO (nil per os - nothing by mouth) for a specified period, typically 6-8 hours, before surgery to prevent aspiration during general anesthesia.
  \item \textbf{Prophylactic Antibiotics:} Intravenous prophylactic antibiotics are administered shortly before the incision to reduce the risk of surgical site infection.
\end{itemize}

\section*{Contraindications \& Precautions}
\textbf{Absolute Contraindications:}
\begin{itemize}
  \item Complete absence of the cochlea or auditory nerve (cochlear aplasia or auditory nerve aplasia), as there is no target for stimulation.
  \item Active infection of the middle ear or mastoid, which must be resolved prior to surgery to prevent device infection.
  \item Active retrocochlear pathology (e.g., untreated acoustic neuroma, active demyelinating disease) that would preclude safe surgery or device function.
  \item Uncontrolled medical conditions (e.g., severe cardiac disease, uncontrolled diabetes, severe coagulopathy) that significantly increase surgical or anesthesia risk.
  \item Profound hearing loss due to central auditory processing disorder without peripheral pathology, as the implant targets the peripheral auditory system.
  \item Unrealistic expectations or unwillingness to participate in the intensive and prolonged post-operative auditory rehabilitation program.
\end{itemize}

\textbf{Relative Contraindications and Precautions:}
\begin{itemize}
  \item Incomplete cochlear ossification or malformations (e.g., Mondini dysplasia, common cavity deformity) may make electrode insertion challenging and require specialized surgical techniques or modified electrode arrays.
  \item Active autoimmune disease affecting the inner ear, which may require immunosuppressive therapy and careful monitoring.
  \item Chronic otitis media with effusion or cholesteatoma, which should be surgically managed and resolved prior to implantation to prevent infection.
  \item Significant cognitive impairment or developmental delay that might hinder rehabilitation, requiring careful consideration, extensive family support, and realistic goal setting.
  \item Certain medical conditions (e.g., severe heart disease, uncontrolled diabetes) that increase surgical risk, necessitating thorough preoperative optimization and multidisciplinary management.
  \item Future need for magnetic resonance imaging (MRI): While most modern implants are MRI-compatible up to certain field strengths (e.g., 1.5T or 3.0T), specific protocols and precautions (e.g., magnet removal, head bandage) are often required, and some older implants may be contraindicated for MRI.
  \item Advanced age: While there is no upper age limit, advanced age with significant comorbidities may increase surgical and anesthesia risk, requiring careful patient selection.
  \item Residual hearing: In some cases, patients with significant residual low-frequency hearing may be candidates for hybrid implants that combine acoustic amplification with cochlear stimulation, aiming to preserve natural hearing.
\end{itemize}

\section*{Complications \& Risks}
As with any surgical procedure, cochlear implant surgery carries inherent risks, which can be broadly categorized as general surgical risks and procedure-specific complications.

\begin{itemize}
  \item \textbf{General Surgical and Anesthesia Risks:}
    \begin{itemize}
      \item \textit{Bleeding:} Hematoma formation at the surgical site, which may rarely require re-exploration.
      \item \textit{Infection:} Wound infection, mastoiditis, or, rarely, meningitis (especially if not vaccinated against \textit{Streptococcus pneumoniae}).
      \item \textit{Anesthesia Complications:} Allergic reactions to medications, respiratory or cardiac complications, nausea, vomiting, or prolonged recovery from anesthesia.
      \item \textit{Pain:} Postoperative pain at the incision site, usually manageable with analgesics.
    \end{itemize}

  \item \textbf{Procedure-Specific Risks:}
    \begin{itemize}
      \item \textit{Facial Nerve Injury:} The facial nerve runs through the mastoid bone; injury can lead to temporary or permanent facial weakness or paralysis. This is a rare but serious complication (incidence 0.1-0.7\%).
      \item \textit{Cerebrospinal Fluid (CSF) Leak:} Leakage of CSF from the inner ear or mastoid, which may require further surgical intervention to repair.
      \item \textit{Tinnitus:} Pre-existing tinnitus may worsen, improve, or new tinnitus may develop.
      \item \textit{Dizziness or Vertigo:} Temporary or persistent imbalance or spinning sensation due to inner ear manipulation, usually resolving within weeks to months.
      \item \textit{Taste Disturbance:} Injury to the chorda tympani nerve, which passes through the middle ear, can cause altered taste sensation (dysgeusia) on the ipsilateral side of the tongue, usually temporary.
      \item \textit{Device Malfunction or Failure:} The internal or external components of the implant may fail over time, requiring revision surgery to replace the device (incidence 1-3\% over 10 years).
      \item \textit{Electrode Misplacement or Migration:} The electrode array may not be fully inserted or may migrate out of the cochlea, leading to suboptimal performance and potentially requiring revision.
      \item \textit{Non-Auditory Stimulation:} Electrical stimulation of adjacent nerves (e.g., facial nerve) can cause twitching, discomfort, or other non-auditory sensations, usually manageable with programming adjustments.
      \item \textit{Lack of Benefit or Dissatisfaction:} Despite successful surgery, some individuals may not achieve the expected level of hearing improvement or may be dissatisfied with the sound quality, requiring extensive rehabilitation and counseling.
      \item \textit{Residual Hearing Loss:} Any remaining natural hearing in the implanted ear is typically lost during the procedure, particularly with full electrode insertion.
      \item \textit{Vestibular Dysfunction:} Damage to the vestibular system during surgery can lead to persistent balance issues.
    \end{itemize}
\end{itemize}

\section*{Postoperative Recovery Timeline}
Immediately after cochlear implant surgery, the patient is transferred to the post-anesthesia care unit (PACU) for close monitoring as they recover from general anesthesia. Pain medication will be administered as needed, and a pressure dressing may be in place over the surgical site to minimize swelling and hematoma. Most patients experience mild discomfort and some swelling behind the ear, which can be managed with oral analgesics. They are typically discharged home within 24 hours, often the same day, with instructions for wound care, pain management, and activity restrictions.

Over the first 1-2 weeks, patients are advised to avoid strenuous activities, heavy lifting, and getting the incision wet. The incision site should be kept clean and dry, and any dressing changes should be performed as instructed. The external speech processor is not worn during this initial healing period. Approximately 3-6 weeks post-surgery, once the incision has fully healed and swelling has subsided, the internal implant is activated, and the external speech processor is fitted and programmed for the first time. This initial programming session, known as "mapping," is a crucial step where the audiologist adjusts the electrical stimulation levels for each electrode to optimize sound perception. This process is iterative, requiring multiple follow-up mapping sessions over several months to fine-tune the device settings and maximize auditory performance. Long-term, recipients engage in extensive auditory rehabilitation and training to learn to interpret the new electrical sounds, which can take months to years to fully adapt to and achieve optimal speech understanding.

\section*{Surgical Findings \& Pathology}
During cochlear implant surgery, the surgeon meticulously documents intraoperative findings related to the anatomy of the temporal bone and inner ear. Key observations include the patency of the cochlea, noting any degree of ossification (bony growth within the cochlea, often seen after meningitis) or malformations (e.g., Mondini dysplasia, common cavity deformity) that might affect electrode insertion. The integrity and course of the facial nerve are carefully noted, especially during the facial recess approach. The condition of the round window membrane and the ease of electrode array insertion into the scala tympani are also critical findings.

Unlike many other surgical procedures, cochlear implant surgery typically does not involve the resection of tissue for histopathological examination. The primary "pathology" in this context is the underlying severe-to-profound sensorineural hearing loss, which is diagnosed preoperatively. The surgical findings focus on the anatomical suitability for implantation and the successful, atraumatic placement of the device. The implant itself is a foreign body, and its secure placement and proper function, confirmed by intraoperative testing, are the main "findings" of the procedure. Any unexpected anatomical variations or difficulties during insertion are documented as they may influence postoperative programming and outcomes.

\section*{Expected Postoperative Course}
A normal and successful postoperative course following cochlear implant surgery involves a smooth recovery from anesthesia, minimal pain at the surgical site, and uneventful wound healing. Patients typically experience mild to moderate discomfort behind the ear for a few days, which is well-controlled with oral pain medication. Swelling and bruising around the incision are common and usually resolve within 1-2 weeks. The incision site should heal cleanly without signs of infection (redness, excessive warmth, pus). Once the incision has fully healed (typically 3-6 weeks post-surgery), the external speech processor is activated. Initially, sounds may seem mechanical or unnatural, but with consistent use and dedicated auditory rehabilitation, the brain gradually adapts to the new electrical signals. Over several months, patients typically experience progressive improvements in sound awareness, speech perception in quiet environments, and eventually, speech understanding in more challenging listening situations. The ultimate goal is to achieve functional communication and an enhanced quality of life.

\section*{Postoperative Care \& Follow-up}
Regular and comprehensive follow-up is crucial for optimizing the performance of a cochlear implant and ensuring long-term success. This involves a multidisciplinary team.

\begin{itemize}
  \item \textbf{Initial Surgical Follow-up:} Typically occurs 1-2 weeks after surgery to check wound healing, remove any non-absorbable sutures or staples, and address any immediate concerns.
  \item \textbf{Device Activation and Mapping:} The first activation and programming (mapping) session takes place 3-6 weeks post-surgery. Subsequent mapping sessions are frequent in the first year (e.g., monthly for the first few months, then quarterly) and then annually or as needed to fine-tune the device settings.
  \item \textbf{Audiometric Evaluations:} Regular audiometry and speech perception testing are performed to monitor hearing progress and guide device programming adjustments.
  \item \textbf{Auditory Rehabilitation:} Ongoing auditory training and speech therapy are essential, especially for children, to develop listening skills, speech understanding, and language development.
  \item \textbf{Processor Upgrades:} External speech processors are periodically upgraded (typically every 5-7 years) to incorporate newer technology and improved sound processing strategies.
  \item \textbf{Device Integrity Checks:} Annual checks of the external processor and internal implant integrity are performed to ensure optimal function.
  \item \textbf{Activity Resumption:} Patients are gradually cleared to resume normal activities, though contact sports may require head protection to prevent trauma to the implant site.
  \item \textbf{Warning Signs:} Patients are educated on warning signs requiring immediate medical attention, such as sudden changes in hearing, persistent pain, excessive swelling, redness, fever, facial weakness, or dizziness.
\end{itemize}
Lifelong follow-up with an audiologist and otolaryngologist is recommended to ensure optimal device performance and address any issues.

\section*{Epidemiology}
Cochlear implantation has become a widely accepted and increasingly common intervention for severe-to-profound sensorineural hearing loss globally. The incidence of severe-to-profound hearing loss varies, but it affects approximately 1-3 per 1,000 live births, with a significant proportion of these children being candidates for cochlear implants. In adults, the prevalence of severe-to-profound hearing loss increases with age, often due to presbycusis, ototoxicity, or other acquired causes. Annually, tens of thousands of cochlear implants are performed worldwide, with numbers steadily rising. While initially more common in children, the adult population now represents a substantial and growing proportion of recipients, reflecting expanded candidacy criteria and improved outcomes. There is no significant sex predilection for the underlying hearing loss or the procedure itself. Common indications include congenital deafness in children and progressive, age-related, or sudden sensorineural hearing loss in adults. Mortality directly attributable to cochlear implant surgery is exceedingly rare, estimated to be less than 0.01\%. Morbidity, while low, includes risks such as facial nerve injury (0.1-0.7\%), infection (1-5\%), and device failure (1-3\% over 10 years), which are generally well-managed.

\section*{Outcomes \& Prognosis}
The outcomes and prognosis for cochlear implant recipients are generally excellent, with high rates of functional success and significant improvements in quality of life. For children with congenital deafness, early implantation (ideally before 12-18 months of age) is consistently associated with the best speech and language development outcomes, often allowing them to achieve age-appropriate communication skills and integrate into mainstream education. Adults with post-lingual deafness typically experience substantial improvements in speech understanding, particularly in quiet environments, with many achieving open-set speech recognition and improved communication abilities. Success rates for achieving functional hearing are high, often exceeding 80-90\% in appropriately selected candidates.

Prognostic factors influencing outcomes include the duration of deafness (shorter duration generally leads to better outcomes), age at implantation (earlier for children, younger for adults with post-lingual loss), etiology of hearing loss, integrity of the auditory nerve, and commitment to post-operative auditory rehabilitation. While the implant does not restore normal hearing, it provides a robust auditory signal that the brain learns to interpret. Device longevity is also favorable, with internal components designed to last many years, though external processors are typically upgraded every 5-7 years to incorporate technological advancements. Recurrence of hearing loss is not applicable as the underlying pathology is bypassed, but device failure or complications can necessitate revision surgery. Overall, cochlear implants offer a highly effective and durable solution for severe-to-profound sensorineural hearing loss, profoundly impacting communication, social integration, and overall well-being.

\section*{References}
\begin{enumerate}
  \item MedlinePlus. Cochlear implants. U.S. National Library of Medicine; 2024. Available from: https://medlineplus.gov/cochlearimplants.html
  \item Flint PW, Haughey BH, Lund VJ, et al. Cummings Otolaryngology: Head \& Neck Surgery. 7th ed. Philadelphia, PA: Elsevier; 2021.
  \item Kliegman RM, St. Geme JW, Blum NJ, Shah SS, Tasker RC, Wilson KM. Nelson Textbook of Pediatrics. 22nd ed. Philadelphia, PA: Elsevier; 2024.
  \item American Academy of Otolaryngology--Head and Neck Surgery. Clinical Consensus Statement: Cochlear Implants. 2023. Available from: https://www.entnet.org/resource/clinical-consensus-statement-cochlear-implants/
\end{enumerate}

\clearpage